%------------------------------------------------------------------------------%
% \chapter{Matrizes}
%------------------------------------------------------------------------------%

%------------------------------------------------------------------------------%
\begin{exercicios}{Determinante de Uma Matriz e Suas Propriedades}
    
    
% 1
\begin{exercicio}
    Calcule os determinantes
    \begin{mathlist}[2]
        \item \left|\begin{array}{rr}
            -3 & -1 \\
            2 &  4
        \end{array} \right|
        \item \left|\begin{array}{rr}
            13 & 7 \\
            11 & 5
        \end{array} \right|
        \item \left|\begin{array}{rr}
            3 & 1 \\
            5 & 2
        \end{array} \right|
        \item \left|\begin{array}{cc}
            \log(a) &     \log(b) \\
            \sfrac{1}{2} & \sfrac{1}{4}
        \end{array} \right|
        \item \left|\begin{array}{cc}
            2m^2 & 2m^4-m \\
            m &  m^3-1
        \end{array} \right|
    \end{mathlist}
    \begin{answers}
        \begin{mathlist}[3]
            \item -10
            \item -12
            \item 1
            \item \log\frac{\sqrt[4]{a}}{\sqrt{b}}
            \item - m^{2}
        \end{mathlist}
    \end{answers}
\end{exercicio}

% 2
\begin{exercicio}
    Calcule os determinantes
    \begin{mathlist}[2]
        \item \left|\begin{array}{rrr}
            1 & 1 & 0 \\
            0 & 1 & 0 \\
            0 & 1 & 1
        \end{array} \right|
        \item \left|\begin{array}{rrr}
            1 & 3 &  2 \\
            -1 & 0 & -2 \\
            2 & 5 &  1
        \end{array} \right|
        \item \left|\begin{array}{rrr}
            -3 & 1 & -7 \\
            2 & 1 &  3 \\
            5 & 4 &  2
        \end{array} \right|
        \item \left|\begin{array}{rrr}
            9 & 7 & 11 \\
            -2 & 1 & 13 \\
            5 & 3 &  6
        \end{array} \right|
        \item \left|\begin{array}{rrr}
            0 & a & c \\
            -c & 0 & b \\
            a & b & 0
        \end{array} \right|
        \item \left|\begin{array}{rrr}
            2 & -1 & 0 \\
            m &  n & 2 \\
            3 &  5 & 4
        \end{array} \right|
    \end{mathlist}
  \begin{answers}
      \begin{mathlist}[3]
        \item 1
        \item -9
        \item 20
        \item 121
        \item b \left(a - c\right) \left(a + c\right)
        \item 4 m + 8 n - 26
    \end{mathlist}
  \end{answers}
\end{exercicio}


% 3
\begin{exercicio}
    Calcule os determinantes das matrizes usando o Teorema de Laplace
    \begin{mathlist}
        \item M = \left|\begin{array}{rrrr}
            3 & 4 &  2 &  1 \\
            5 & 0 & -1 & -2 \\
            0 & 0 &  4 &  0 \\
            -1 & 0 &  3 &  3
        \end{array} \right|
        \item M = \left|\begin{array}{rrrr}
            0 & a & b & 1 \\
            0 & 1 & 0 & 0 \\
            a & a & 0 & b \\
            1 & b & a & 0
        \end{array} \right|
        \item M = \left|\begin{array}{rrrr}
            1 & 3 & 2 & 0 \\
            3 & 1 & 0 & 2 \\
            2 & 3 & 0 & 1 \\
            0 & 2 & 1 & 3
        \end{array} \right|
        \item M = \left|\begin{array}{rrrrr}
            1 & 2 &  3 & 4 & 5 \\
            0 & a & -1 & 3 & 1 \\
            0 & 0 &  b & 2 & 3 \\
            0 & 0 &  0 & c & 2 \\
            0 & 0 &  0 & 0 & d
        \end{array} \right|
        \item M = \left|\begin{array}{cccccc}
            x & 0 & 0 & 0 & 0 & 0 \\ 
            a & y & 0 & 0 & 0 & 0 \\ 
            l & p & z & 0 & 0 & 0 \\ 
            m & n & p & x & 0 & 0 \\ 
            b & c & d & e & y & 0 \\ 
            a & b & c & d & e & z
        \end{array} \right|
    \end{mathlist}        
  \begin{answers}
      \begin{mathlist}[3]
        \item -208
        \item a^{2} + b^{2}
        \item 48
        \item a b c d
        \item x^{2} y^{2} z^{2}
    \end{mathlist}
  \end{answers}
\end{exercicio}

\begin{exercicio}
    Calcule o valor do determinante
    \( \left|\begin{array}{cccc}
    	1 & a & a & a \\
    	1 & 0 & a & a \\
    	1 & 1 & 0 & a \\
    	1 & 1 & 1 & 1
    \end{array} \right| \)
    \begin{answers}
        \(\qquad a^{2} \left( 1 - a \right) \)
    \end{answers}
\end{exercicio}

\begin{exercicio}
    Calcule o valor de $A$ sabendo que
    \[A = \left|\begin{array}{cccc}
    	1 & 2 & 4  &  8  \\
    	1 & 3 & 9  & 27  \\
    	1 & 4 & 16 & 64  \\
    	1 & 5 & 25 & 125
    \end{array} \right|\]
    \begin{answers}
          \(\qquad A = 12\)
    \end{answers}

\end{exercicio}

    
    \begin{exercicio}
        Determine $x$ tal que
        \begin{mathlist}
            \item \left|\begin{array}{cc}
            	2x & 3x+2 \\
            	 1 &    x
            \end{array} \right| = 0
            \item \left|\begin{array}{cc}
            	 2x  & x-2  \\
            	4x+5 & 3x-1
            \end{array} \right| = 11
        \end{mathlist}
        \begin{answers}
            \begin{mathlist}[2]
                \item x_1 = - \sfrac{1}{2} \quad x_2 = 2
                \item x_1 = -1 \quad x_2 = \sfrac{1}{2}
            \end{mathlist}
        \end{answers}
    \end{exercicio}
    
    \begin{exercicio}
        Determine $x$ tal que
        \begin{mathlist}
            \item \left|\begin{array}{rcr}
            	1 &  x  & x \\
            	2 & 2x  & 1 \\
            	3 & x+1 & 1
            \end{array} \right| = 0
            \item \left|\begin{array}{rrr}
            	1 &  x & 1 \\
            	1 & -1 & x \\
            	1 & -x & 1
            \end{array} \right| = 0
            \item \left|\begin{array}{rrr}
            	 1 &  x &  2 \\
            	-2 &  x & -4 \\
            	 1 & -3 & -x
            \end{array} \right| = 0
            \item \left|\begin{array}{ccr}
            	x-1 &  2  &  x \\
            	 0  &  1  & -1 \\
            	3x  & x+1 & 2x
            \end{array} \right| = 
            \left|\begin{array}{cc}
            	3x & 2x \\
            	4  & -x
            \end{array} \right|
        \end{mathlist}
        \begin{answers}
            \begin{mathlist}[2]
                \item x_1 = \sfrac{1}{2}
                \item x_1 = 0 \quad x_2 = 1
                \item x_1 = -2 \quad x_2 = 0
                \item x_1 = - \sfrac{\sqrt{3}}{3} \quad x_2 = \sfrac{\sqrt{3}}{3}
            \end{mathlist}
        \end{answers}
    \end{exercicio}

    \begin{exercicio}
        Encontre os valores de $a$ para os quais a matriz $A$ possui inversa.
        \[A = \left[ 
            \begin{array}{rrr}
            	1 & 1 & 0 \\
            	1 & 0 & 0 \\
            	1 & 2 & a
            \end{array}
        \right] \]
        \begin{answers} \(\qquad a\neq 0 \)
        \end{answers}
    \end{exercicio}
    
    \begin{exercicio}
        Calcule o determinante da matriz $M$ e determine para quais valores de $x$ ela possui inversa
        \[M = \left[ 
            \begin{array}{cccc}
            	4 & x & 0 & x \\
            	3 & 0 & 2 & 0 \\
            	0 & 2 & 0 & 0 \\
            	x & 0 & 1 & 0
            \end{array}
        \right]\]
        \begin{answers} \(\qquad x\neq 0 \) e \( x\neq \sfrac{3}{2} \)
        \end{answers}
    \end{exercicio}

    \begin{exercicio}
        Calcule o determinante da matriz 
        \[ A = \left[\begin{array}{cc}
            	e^{rt}  &   te^{rt}    \\
            	re^{rt} & (1+rt)e^{rt}
            \end{array} \right] 
        \]
        \begin{answers} \(\qquad e^{2 r t}\)
        \end{answers}
    \end{exercicio}
    
    \begin{exercicio}
        Sabendo que $\det A = -3$ calcule
        \begin{mathlist}[3]
            \item \det A^2
            \item \det A^3
            \item \det A^t
        \end{mathlist}
        \begin{answers}
            \begin{mathlist}[3]
                \item 9
                \item -27
                \item(-3)^t
            \end{mathlist}
        \end{answers}
    \end{exercicio}
     
     \begin{exercicio}
         Determine os valores de $\lambda$ tais que \[\det(A-\lambda I) = 0\] para as matrizes
         \begin{mathlist}
            \item A = \left[
                          \begin{array}{ccc}
                              0 & 1 & 2 \\ 
                              0 & 0 & 3 \\ 
                              0 & 0 & 0
                          \end{array}   
                      \right]
            \item A = \left[
                        \begin{array}{rrr}
                             1 & 0 &  0 \\ 
                            -1 & 3 &  0 \\ 
                             3 & 2 & -2
                        \end{array}   
                      \right]
        \end{mathlist}
    \end{exercicio}
     
    \begin{exercicio}
        Calcule o determinantes das matrizes
        \begin{mathlist}[2]
            \item A = \left[
                        \begin{array}{cc}
                            4 & 2 \\ 
                            3 & 2
                        \end{array} 
                      \right]
            \item A = \left[
                        \begin{array}{cc}
                          1 & 2 \\ 
                          0 & 4
                        \end{array} 
                      \right]
            \item AB
            \item 2A
            \item A^tB
            \item A^2B^{-1}
        \end{mathlist}
    \end{exercicio}

    \begin{exercicio}
        Calcule o determinante da matriz 
        \( \left[\begin{array}{cc}
        	a & b \\
        	b & a
        \end{array} \right] \), 
        sabendo que $2a = e^x + e^{-x}$ e $2b = e^x - e ^{-x}$.
    \end{exercicio}
    
    \begin{exercicio}
        Seja \( A = \left[\begin{array}{ccc}
        	2 & n & 2 \\
        	5 & 1 & m \\
        	7 & 3 & 7
        \end{array} \right] \)
        \begin{alphalist}
            \item Ache um valor de $m$, tal que $\det A = 0$, qualquer que seja $n$.
            \item Esse valor de $m$ é o único com essa propriedade?
        \end{alphalist}
    \end{exercicio}
    
    \begin{exercicio}
        A matriz real $A = (a_{ij})$ quadrada e de ordem~3, é definida pela regra
        \[a_{ij} = \begin{cases}
        	0,            & \text{se } i > j \\
        	i^j,          & \text{se } i = j \\
        	(-1)^{i + 1}, & \text{se } i < j
        \end{cases}\]
        \begin{alphalist}
            \item Escreva a matriz $A$ na forma usual de linhas e colunas, onde o valor numérico de $a_{ij}$ está na $i$-ésima linha e na $j$-ésima coluna de $A$.
            \item Calcule o determinante $D$ da matriz $A$.
        \end{alphalist}
    \end{exercicio}
    
    \begin{exercicio}
        Determine o valor de $x$ para que o determinante da matriz $C = AB^t$ seja igual a $602$, onde
        \[A = \left[\begin{array}{rrr}
        	1 & 2 & -3 \\
        	4 & 1 &  2
        \end{array} \right]
        \qquad
        B = \left[\begin{array}{crr}
        	x - 1 & 8 & -5 \\
        	 -2   & 7 &  4
        \end{array} \right]\]
        e $B^t$ é a matriz transposta de $B$.
    \end{exercicio}
    
    \begin{exercicio}
        Duas matrizes de ordem $n$ são inversas quando o produto delas é a identidade de ordem $n$. A inversa da matriz
        \[M = \left[\begin{array}{rrr}
        	1 &  1 & 2 \\
        	1 & -1 & 3 \\
        	2 &  3 & 1
        \end{array} \right]\]
        é a matriz
        \[
        M^{-1} = \left[\begin{array}{ccc}
        	x & 1 & 1 \\
        	1 & y & z \\
        	1 & z & w
        \end{array} \right]\]
        Encontre
        \begin{alphalist}
            \item o valor da expressão $A = x - 3y - z + w$.
            \item $\det(2M)$.
        \end{alphalist}
    \end{exercicio}
    
    \begin{exercicio}
        Determine o conjunto solução da equação
        \[\left|\begin{array}{lrc}
        	3^x &  1 & 0 \\
        	4   &  2 & 2 \\
        	2   & -1 & 3
        \end{array} \right| = \left|\begin{array}{cc}
            4 & 2 \\ 
            2 & 5
        \end{array} \right|\]
        \begin{answers} \( \qquad x = 1\)
        \end{answers}
    \end{exercicio}
    
    \begin{exercicio}
        Determine $x$ de modo que
        \[\left|\begin{array}{crl}
        	1 &  1 & 1   \\
        	2 & -3 & x   \\
        	4 &  9 & x^2
        \end{array} \right| \geq 0 \]
    \end{exercicio}
    
    \begin{exercicio}
        Resolva a equação
        \[\left|\begin{array}{cccc}
            1 & 1 & 0 & 1 \\ 
            x & 1 & x & 0 \\ 
            x & x & 1 & 0 \\ 
            0 & 1 & 0 & 1
        \end{array} \right| = 0\]
        \begin{answers} \(\quad x_1 = -1 \quad x_2 = 1 \)
        \end{answers}
    \end{exercicio}
    
    \begin{exercicio}
        Desenvolva o determinante a seguir
        \[x = \left|\begin{array}{rrrr}
        	a & b  &  c & d \\
        	0 & 1  &  2 & 0 \\
        	1 & -1 &  0 & 1 \\
        	0 & 1  & -1 & 0
        \end{array} \right|\]
    \end{exercicio}
    
    \begin{exercicio}
        Sabendo que
        \( \left|\begin{array}{ccc}
        	a & x & 1 \\
        	b & y & 2 \\
        	c & z & 3
        \end{array} \right| = 4  \) \\
        calcule
        \( M = \left|\begin{array}{ccc}
        	2a & 3x & -1 \\
        	2c & 3z & -3 \\
        	2b & 3y & -2
        \end{array} \right|  \)
    \end{exercicio}
    
    \begin{exercicio}
        Sem calcular os determinantes a seguir, justifique a sua nulidade
        \begin{mathlist}[2]
            \item \left|\begin{array}{rrr}
            	2 & 0 & -1 \\
            	5 & 0 &  5 \\
            	4 & 0 &  2
            \end{array} \right|
            \item \left|\begin{array}{rrr}
            	 a &  b &  c \\
            	ka & kb & kc \\
            	 x &  y &  m
            \end{array} \right|
            \item \left|\begin{array}{rrr}
            	 5 &  4 & 3 \\
            	-4 & -2 & 2 \\
            	 4 &  5 & 6
            \end{array} \right|
            \item \left|\begin{array}{rrr}
            	1 & 2 & 5 \\
            	m & 5 & 4 \\
            	1 & 2 & 5
            \end{array} \right|
            \item \left|\begin{array}{rrr}
            	2 & 3 &  4 \\
            	1 & 2 &  7 \\
            	3 & 5 & 11
            \end{array} \right|
            \item \left|\begin{array}{rrrr}
            	 1 &  2 &  3 &  4 \\
            	 5 &  6 &  7 &  8 \\
            	 9 & 10 & 11 & 12 \\
            	13 & 14 & 15 & 16
            \end{array} \right|
        \end{mathlist}
    \end{exercicio}
    
    \begin{exercicio}
        Calcule $x$ para que os determinantes das matrizes sejam iguais 
        \[A=\left[\begin{array}{crcc}
        	x + 1 &  0 & 0 & 2 \\
        	  0   & -2 & 0 & 5 \\
        	  x   &  0 & 3 & 0 \\
        	  0   &  1 & 0 & 4
        \end{array} \right] 
        \qquad
        B=
        \left[\begin{array}{cc}
        	1  & x \\
        	78 & 0
        \end{array} \right]\]
    \end{exercicio}
    
    \begin{exercicio}
        Determine as condições que $x$ deve satisfazer para que a matriz $A$ seja inversível
        \[A = \left[\begin{array}{cccc}
            1 & 2 & 3 & 4 \\ 
            1 & 3 & x & 5 \\ 
            1 & 3 & 4 & 3 \\ 
            1 & 6 & 5 & x
        \end{array} \right]\]
    \end{exercicio}
    
    \begin{exercicio}
        Seja a matriz 
        \[A = \left[\begin{array}{ccc}
            k & k & k \\ 
            k & k & 5 \\ 
            k & 5 & 5
        \end{array} \right]\]
        Determine $k \in \R$ de modo que a matriz $A$ não admita inversa.
    \end{exercicio}
    
    \begin{exercicio}
        Dadas as matrizes
        \[A = \left[\begin{array}{cc}
        	1 & 2 \\
        	2 & 0
        \end{array} \right]
        \qquad
        B = \left[\begin{array}{cr}
        	4 & -4 \\
        	1 &  4
        \end{array} \right]\]
        calcule o determinante da matriz $(AB)^{-1}$.
    \end{exercicio}
     
\end{exercicios} % Determinante de Uma Matriz e Suas Propriedades

%------------------------------------------------------------------------------%
